\section{Spatial rank of the sampled normalized power}
\label{sec:lvv_sampled_power_rank}

Consider a chain of $L\ge0$ sites, each of dimension $q\ge1$.
Let $I$ index its adjacent pairs. An operator $P_i$ is supported on pair
$i,i+1$ for every $i\in I$. Adjoin a letter $\ast$, put $P_\ast=1$, and
write $P_w=P_{w(0)}\cdots P_{w(m-1)}$ for a word $w$ of length $m$.
Let $C_a(w)$ be the number of occurrences of the letter $a$ in $w$.
Denote operator Schmidt rank across the cut after $k$ sites by
$\operatorname{SR}_k$.

\begin{theorem}[The unique crossing letter]
    \label{thm:lvv_padded_nn_word_rank}
    \lean{TNLean.MPS.AreaLaw.operatorSchmidtRank_padded_nn_word_le}
    \leanok
    At an interior cut $0<k<L$, let $a_k$ be the letter on the adjacent
    pair crossing that cut. Then
    \begin{align}
        \operatorname{SR}_k(P_w)\le q^{2C_{a_k}(w)}.
        \label{eq:lvv_padded_nn_word_rank}
    \end{align}
    This is the ordered-word estimate of~\cite[Section~2.3.1]{Landau2015Polynomial}, with the paired-index correction
    recorded in~\cite{gap:lvv13_paired_matrix_decomposition}.
\end{theorem}
\begin{proof}\leanok
    A letter away from the crossing pair acts on one side and contributes
    rank at most one. The crossing pair meets either side in one site,
    so its rank is at most $q^2$. The general ordered-support product bound
    therefore gives~(\ref{eq:lvv_padded_nn_word_rank}), with the actual
    multiplicity of the crossing letter. No commutation is required.
\end{proof}

\begin{theorem}[One sample family, all spatial cuts]
    \label{thm:lvv_padded_nn_sample_rank}
    \lean{TNLean.MPS.AreaLaw.operatorSchmidtRank_padded_nn_word_le_of_counts,
        TNLean.MPS.AreaLaw.operatorSchmidtRank_padded_nn_sample_sum_le,
        TNLean.MPS.AreaLaw.operatorSchmidtRank_sampled_nn_complements_le}
    \leanok
    \uses{thm:lvv_padded_nn_word_rank}
    Let $\sigma$ be a family of $\ell$ words such that
    $C_a(\sigma(b))\le u$ for every sample $b$ and letter $a$.
    For arbitrary complex coefficients $z_b$ and every cut $k\ge0$,
    \begin{align}
        \operatorname{SR}_k(P_{\sigma(b)})
        &\le q^{2u},\\
        \operatorname{SR}_k\left(\sum_b z_bP_{\sigma(b)}\right)
        &\le \ell q^{2u}.
        \label{eq:lvv_padded_nn_sample_rank}
    \end{align}
    The same bound holds after scalar normalization, and applies in
    particular to the complements $P_i=1-h_i$ of adjacent-pair terms.
\end{theorem}
\begin{proof}\leanok
    \uses{thm:lvv_padded_nn_word_rank}
    At an interior cut, apply~(\ref{eq:lvv_padded_nn_word_rank}) to each word.
    At an endpoint cut, every letter acts on one side, so every word has
    rank at most one. Finally use subadditivity and the fact that scalar
    multiplication does not increase operator Schmidt rank. The identity
    and $h_i$ are supported on the same pair, hence so is $1-h_i$.
\end{proof}

\begin{theorem}[The normalized power sample with spatial rank control]
    \label{thm:lvv_sampled_nn_power_rank}
    \lean{TNLean.MPS.AreaLaw.exists_sampled_nn_power_approximation_with_rank}
    \leanok
    \uses{thm:lvv_padded_nn_sample_rank}
    Let $0\le h_i\le1$ be supported on the adjacent pairs, and set
    $H=\sum_{i\in I}h_i$, $n=|I|+1$. Suppose that $\Omega$ is a unit vector,
    $H\Omega=E_0\Omega$, and
    $H-E_0\,1\ge\Delta(1-\Pi)$, where $\Delta>0$ and
    $\Pi=|\Omega\rangle\langle\Omega|$.
    Choose $Q\ge1$, $\eta>0$, and an integer $p\ge1$ with $2q^L<4^p$.
    Put
    \begin{align}
        C&=(1-E_0/n)^{-1},\\
        m&=\left\lceil\frac{n-E_0}{\Delta}\log Q\right\rceil,\\
        S&=n^2\exp\left(\frac{2E_0\log Q}{\Delta}\right),\\
        \ell&=\left\lceil\frac{16p^2S}{\eta^2}\right\rceil+1,\\
        u&=\left\lceil\frac{m/n+\log(2n\ell)}{\log2}\right\rceil.
        \label{eq:lvv_sampled_nn_power_parameters}
    \end{align}
    Then $\ell>0$, and there is one family $\sigma$ of $\ell$ words of length
    $m$ over $I\sqcup\{\ast\}$ such that, for
    \begin{align}
        A=[C(1-H/n)]^m,\qquad
        K=\frac{C^m}{\ell}\sum_bP_{\sigma(b)},
        \label{eq:lvv_sampled_nn_power_operator}
    \end{align}
    one has
    \begin{align}
        \|K-A\|&<\eta,\\
        \|K-\Pi\|&<\eta+Q^{-1},\\
        C_a(\sigma(b))&<u\quad\text{for every }b,a,\\
        \operatorname{SR}_k(K)&\le\ell q^{2u}\quad\text{for every }k\ge0.
        \label{eq:lvv_sampled_nn_power_conclusions}
    \end{align}
    The ground-energy dependence agrees with the published sampling
    argument~\cite[Methods, Proposition~1, p.~6]{Landau2015Published}.
\end{theorem}
\begin{proof}\leanok
    \uses{thm:lvv_padded_nn_sample_rank}
    The finite uniform sampling theorem gives the first three conclusions
    on a single actual word family. Its proof derives the word mean from
    the expansion of $[C(1-H/n)]^m$ and uses the norm bound $\|P_w\|\le1$.
    The displayed choices make the matrix failure probability strictly
    below $1/2$ and the frequency failure probability at most $1/2$.
    Apply Theorem~\ref{thm:lvv_padded_nn_sample_rank} to this same family,
    with coefficients $C^m/\ell$.
\end{proof}

The factor $S=n^2Q^{2E_0/\Delta}$ remains in the sample count. Thus this
construction does not establish the unrestricted polynomial sample count
printed in the first preprint version. The distinction between that
version and the published theorem is explained in the note on ground-energy
normalization accompanying this development.

\begin{theorem}[Sampled left factors for a specified ground energy]
    \label{thm:lvv_sampled_left_error_reduction}
    \lean{TNLean.MPS.AreaLaw.exists_sampled_nn_errorReduction}
    \leanok
    \uses{thm:lvv_sampled_nn_power_rank,thm:lvv_simultaneous_cut_decomposition,
        thm:lvv_projection_error_reduction,thm:lvv_operator_bond_growth}
    Under the physical hypotheses of Theorem~\ref{thm:lvv_sampled_nn_power_rank},
    let $0<\delta\le1$ and $p\ge1$ satisfy $2q^L<4^p$.
    In~(\ref{eq:lvv_sampled_nn_power_parameters}), choose
    $Q=8/\delta$, $\eta=\delta/8$, and put $r=\ell q^{2u}$.
    The actual ground energy $E_0$ is specified before these parameters
    and the sampled family are chosen.

    There is one sampled operator $K$, with its actual word family as
    in~(\ref{eq:lvv_sampled_nn_power_operator}), such that
    $\|K-\Pi\|<\delta/4$ and the following holds at every cut $i\le L$.
    There are $a\le r$ pairs of operators $A_j,B_j$, acting on the first
    $i$ sites and the remaining sites respectively, such that
    \begin{align}
        K=\sum_{j=1}^{a}A_j\otimes B_j.
        \label{eq:lvv_sampled_left_factor_decomposition}
    \end{align}
    Every $A_j$ has operator Schmidt rank at most $r$ at each internal
    cut of the first $i$ sites. These factors may be chosen so that for
    every $(i,s,b,1/2)$-viable finite set $S$, the actual family
    \begin{align}
        S'=\{A_js:1\le j\le a,\ s\in S\}
        \label{eq:lvv_sampled_left_factor_family}
    \end{align}
    is $(i,as,r\max(1,b),\delta)$-viable.
    In particular, $|S'|\le as\le rs$.
\end{theorem}
\begin{proof}\leanok
    \uses{thm:lvv_sampled_nn_power_rank,thm:lvv_simultaneous_cut_decomposition,
        thm:lvv_projection_error_reduction,thm:lvv_operator_bond_growth}
    The chosen parameters give
    $\eta+Q^{-1}=\delta/4<1/2$.
    Theorem~\ref{thm:lvv_sampled_nn_power_rank} supplies a single actual
    sampled family with this projection error and rank at most $r$ at
    all cuts. Theorem~\ref{thm:lvv_simultaneous_cut_decomposition}
    constructs~(\ref{eq:lvv_sampled_left_factor_decomposition}), with
    $a\le r$ and the asserted internal ranks of its left factors.
    These choices precede the choice of $S$.

    The deterministic projection estimate,
    Theorem~\ref{thm:lvv_projection_error_reduction}, gives overlap loss
    at most $4(\delta/4)=\delta$.
    Theorem~\ref{thm:lvv_operator_bond_growth} bounds the bonds of
    $A_js$ by $r\max(1,b)$, and there are at most $a|S|$ such vectors.
\end{proof}

The sample count still contains
$\exp(2E_0\log(8/\delta)/\Delta)$. This theorem does not give a construction
uniform in an unknown ground energy or an energy-independent polynomial
sample count. Choosing the word length and relative-error budget without
an exact energy input remains separate; an unknown overall scalar is
immaterial after output normalization.

\begin{theorem}[One sampled family for every unit ground vector]
    \label{thm:lvv_sampled_ground_line_error_reduction}
    \lean{TNLean.MPS.AreaLaw.exists_sampled_nn_errorReduction_forall_ground_vectors}
    \leanok
    \uses{thm:lvv_sampled_left_error_reduction}
    Retain the interaction, gap, accuracy and sampling parameters of
    Theorem~\ref{thm:lvv_sampled_left_error_reduction}, with the actual energy
    $E_0$ specified before sampling. Suppose that there exists a unit vector
    $\Omega_0$ satisfying
    $H\Omega_0=E_0\Omega_0$ and
    $H-E_0\,1\ge\Delta(1-|\Omega_0\rangle\langle\Omega_0|)$.

    One sample family, its operator $K$, and one family of cut factors at
    each cut may be chosen so that all conclusions of
    Theorem~\ref{thm:lvv_sampled_left_error_reduction} hold simultaneously
    for every normalized vector $\Omega$ with $H\Omega=E_0\Omega$.
    Thus neither the sample nor its factors depend on the choice of a unit
    representative of the ground line. The factors also precede the choice
    of every half-viable input family.
\end{theorem}
\begin{proof}\leanok
    \uses{thm:lvv_sampled_left_error_reduction}
    Choose $\Omega_0$ internally and apply the preceding construction once.
    The strictly positive lower gap implies that every vector at energy
    $E_0$ equals its projection onto $\Omega_0$:
    \begin{align}
        \Omega=\langle\Omega_0,\Omega\rangle\Omega_0.
        \label{eq:lvv_sampled_ground_line_uniqueness}
    \end{align}
    If $\Omega$ is normalized, the scalar in this identity has modulus one.
    Consequently
    $|\Omega\rangle\langle\Omega|=|\Omega_0\rangle\langle\Omega_0|$,
    so the same sampled operator has projection error below $\delta/4$
    for every such $\Omega$.

    The cut decomposition and the internal ranks of its left factors depend
    only on $K$. Apply the deterministic error-reduction estimate to those
    same factors and each normalized ground vector. This gives the same
    cardinality, MPS bond and error bounds for every corresponding viable
    input family.
\end{proof}

The sample count retains its dependence on the specified energy $E_0$.
The simultaneous conclusion for all unit ground vectors does not select
the word length or relative-error sample budget without that energy input.
An overall nonzero scalar is immaterial after normalizing the output vectors.

\begin{theorem}[Sampling without scalar normalization from a known energy interval]
    \label{thm:lvv_interval_power_sampling}
    \lean{TNLean.MPS.AreaLaw.exists_interval_sampled_nn_power}
    \leanok
    \uses{thm:lvv_padded_nn_sample_rank,thm:lvv_simultaneous_cut_decomposition,
        thm:lvv_scaled_projection_error_reduction}
    Let $H=\sum_{i\in I}h_i$ be the sum of positive contractions on the
    adjacent pairs of a chain of $L\ge0$ sites of dimension $q\ge1$.
    Put $n=|I|+1$, and fix $0\le a\le n-1$, $\Delta>0$, $Q\ge1$, $\eta>0$,
    and an integer $p\ge1$ such that $2q^L<4^p$.
    Define
    \begin{align}
        E_{\max}&=\min(a+1,n-1), & C_{\max}&=(1-E_{\max}/n)^{-1},\\
        m&=\left\lceil\frac{n-a}{\Delta}\log Q\right\rceil,
        & S&=n^2Q^{4(a+1)/\Delta},\\
        \ell&=\left\lceil\frac{16p^2S}{\eta^2}\right\rceil+1,
        & u&=\left\lceil\frac{m/n+\log(2n\ell)}{\log2}\right\rceil,\\
        r&=\ell q^{2u}.
        \label{eq:lvv_interval_sampling_parameters}
    \end{align}
    There is a family $\sigma$ of $\ell>0$ words of length $m$ such that
    its unscaled sample
    \begin{align}
        T=\frac1\ell\sum_bP_{\sigma(b)}
        \label{eq:lvv_interval_unscaled_sample}
    \end{align}
    satisfies
    \begin{align}
        \|T-(1-H/n)^m\|&<\frac{\eta}{C_{\max}^m},\\
        C_c(\sigma(b))&<u\quad\text{for every }b,c,\\
        \operatorname{SR}_k(T)&\le r\quad\text{at every cut }k.
        \label{eq:lvv_interval_mean_counts_rank}
    \end{align}

    This one family works for every later choice of $E_0\in[a,a+1]$ and
    unit vector $\Omega$ satisfying
    $H\Omega=E_0\Omega$ and $H-E_0\,1\ge\Delta(1-\Pi)$, where
    $\Pi=|\Omega\rangle\langle\Omega|$.
    More precisely, with $\beta=(1-E_0/n)^m$ one has
    \begin{align}
        \beta>0,\qquad \|T-\beta\Pi\|<\beta(\eta+Q^{-1}).
        \label{eq:lvv_interval_relative_projection}
    \end{align}
    At each cut $i\le L$, one may choose a decomposition
    $T=\sum_{j=1}^{t}A_j\otimes B_j$ with $t\le r$ and all internal
    cut ranks of $A_j$ at most $r$, before choosing those ground data.
    If $\eta+Q^{-1}<1/2$, these same left factors send every
    $(i,s,b,1/2)$-viable family to an
    $(i,ts,r\max(1,b),4(\eta+Q^{-1}))$-viable family.
\end{theorem}
\begin{proof}\leanok
    \uses{thm:lvv_padded_nn_sample_rank,thm:lvv_simultaneous_cut_decomposition,
        thm:lvv_scaled_projection_error_reduction}
    The interval estimate gives
    $C_{\max}^{2m}\le n^2Q^{4(a+1)/\Delta}=S$.
    Apply the proved joint finite sampling bound to the actual unscaled
    complement words, at tolerance $\eta/C_{\max}^m$. The chosen count
    and frequency threshold make the sum of the failure bounds less
    than one. This constructs a single family satisfying the first two
    conclusions in~(\ref{eq:lvv_interval_mean_counts_rank}).
    Theorem~\ref{thm:lvv_padded_nn_sample_rank} gives its simultaneous
    spatial ranks; Theorem~\ref{thm:lvv_simultaneous_cut_decomposition}
    constructs the cut factors before any ground data.

    A later unit ground eigenvector gives $E_0\le n-1$ by evaluating
    the original contraction sum. Hence $E_0\le E_{\max}$ and
    $C_{\max}^{-m}\le\beta$. The bracket also gives
    $m\ge(n-E_0)\log Q/\Delta$.
    The existing lower-gap power estimate, rescaled by $\beta$, yields
    $\|(1-H/n)^m-\beta\Pi\|\le\beta/Q$.
    Combining this with the actual sample-mean error proves
    (\ref{eq:lvv_interval_relative_projection}).
    Finally, Theorem~\ref{thm:lvv_scaled_projection_error_reduction}
    divides the perturbation error by the positive amplitude $\beta$,
    giving the asserted error, cardinality and MPS bond bounds.
\end{proof}

The interval $[a,a+1]$ is supplied in advance. This theorem proves the
existence of a finite sample with the stated properties. It neither
selects the interval nor instantiates the source's high success probability
or a complete trial-energy algorithm. The budget retains its dependence
on $a+1$.
